% ============================================================================
\section{Giới thiệu}
\label{sec:intro}
% ============================================================================

\subsection{Đặt vấn đề}

Khối lượng văn bản mà một người cần đọc để tìm một thông tin cụ thể --- giáo trình,
quy chế, tài liệu kỹ thuật, bài viết bách khoa --- tăng nhanh hơn nhiều so với thời
gian họ có. Công cụ tìm kiếm truyền thống giải quyết một nửa bài toán: nó trả về
\emph{tài liệu} có liên quan, nhưng người dùng vẫn phải tự đọc để tìm \emph{câu trả
lời}. \textbf{Đọc hiểu máy} (Machine Reading Comprehension --- MRC) giải quyết nửa còn
lại: cho một đoạn văn và một câu hỏi, hệ thống chỉ thẳng vào đoạn chữ trả lời câu hỏi
đó. Đây là thành phần lõi của trợ lý học tập, hệ thống hỏi đáp trên tài liệu nội bộ và
các pipeline tìm kiếm tri thức hiện đại, nơi một bộ truy hồi chọn đoạn văn và một bộ
đọc hiểu trích xuất câu trả lời.

Với tiếng Anh, bài toán đã được nghiên cứu sâu nhờ các bộ dữ liệu lớn như SQuAD
\cite{rajpurkar2016squad,rajpurkar2018know}. Với tiếng Việt, các bộ dữ liệu UIT-ViQuAD
\cite{nguyen2020viquad} và UIT-ViQuAD~2.0 \cite{nguyen2022vlsp} của Trường Đại học Công
nghệ Thông tin đã tạo nền tảng đánh giá chung. Tuy vậy, tiếng Việt đặt ra những khó khăn
riêng mà phương pháp và công cụ thiết kế cho tiếng Anh không tự động xử lý đúng:

\begin{itemize}
  \item \textbf{Đơn vị chữ viết là âm tiết, không phải từ.} ``Hà Nội'', ``thủ đô'' là
        từ gồm hai âm tiết cách nhau bởi khoảng trắng. Mọi quyết định tách token --- trong
        metric lẫn trong mô hình --- đều ảnh hưởng đến kết quả.
  \item \textbf{Dấu thanh mang nghĩa.} ``hoà'' và ``hoa'' là hai từ khác nhau. Một
        bước chuẩn hoá hoặc giải mã làm mất dấu sẽ âm thầm làm sai kết quả.
  \item \textbf{Không có mạo từ.} Quy tắc chuẩn hoá của SQuAD loại bỏ \emph{a/an/the};
        trong tiếng Việt, những từ gần tương đương như ``các'', ``những'', ``con'' là
        loại từ và có thể là một phần của đáp án đúng.
  \item \textbf{Câu hỏi không có đáp án.} UIT-ViQuAD~2.0 có khoảng một phần ba câu hỏi
        mà đoạn văn \emph{không} trả lời được. Hệ thống phải biết từ chối, không chỉ biết
        tìm.
\end{itemize}

Ngoài khía cạnh ngôn ngữ, đề tài còn đặt ra một thách thức về \emph{phương pháp}. Hướng
dẫn chung của môn học nhấn mạnh: \emph{``Kết quả cao trên benchmark không đồng nghĩa với
khả năng triển khai thực tế. Sinh viên phải nêu rõ giới hạn dữ liệu và tránh khẳng định
vượt quá phạm vi thực nghiệm.''} Vì vậy đồ án không đặt mục tiêu ``đạt số cao nhất'', mà
đặt mục tiêu \emph{mọi con số được trình bày đều đúng, truy vết được và được diễn giải
đúng phạm vi}.

\subsection{Phát biểu bài toán}
\label{sec:problem}

Cho một đoạn văn $C = (c_1, \dots, c_{|C|})$ là chuỗi ký tự và một câu hỏi $Q$. Bài toán
\textbf{đọc hiểu trích xuất} yêu cầu tìm cặp chỉ số ký tự $(s, e)$ với
$0 \le s < e \le |C|$ sao cho đoạn con $C[s{:}e]$ trả lời $Q$; hoặc trả về chuỗi rỗng
$\varnothing$ nếu $C$ không chứa câu trả lời.

\begin{table}[htbp]
\centering
\caption{Đặc tả đầu vào, đầu ra của hệ thống}
\label{tab:io}
\renewcommand{\arraystretch}{1.2}
\begin{tabularx}{\textwidth}{@{}l L@{}}
\toprule
\textbf{Thành phần} & \textbf{Mô tả} \\
\midrule
Đầu vào & \code{context}: đoạn văn Wikipedia tiếng Việt; \code{question}: câu hỏi tự nhiên \\
Đầu ra  & \code{answer}: một chuỗi \emph{con} của \code{context}, hoặc chuỗi rỗng khi không có đáp án \\
Loại bài toán & Trích xuất (extractive), \emph{không} sinh văn bản (generative) \\
Bộ dữ liệu & UIT-ViQuAD~2.0, định dạng SQuAD~2.0 \\
Thang đo & Exact Match (EM) và F1 trên token (âm tiết) \\
\bottomrule
\end{tabularx}
\end{table}

Tính chất \emph{trích xuất} không phải chi tiết phụ: nó cho phép phát biểu một
\textbf{bất biến kiểm tra được} --- với mọi mô hình và mọi đầu vào, đầu ra phải là chuỗi
con của \code{context}. Bất biến này được kiểm thử tự động cho các mô hình và bắt được
cả lỗi lệch offset của tokenizer lẫn trường hợp mô hình ``bịa'' ra chữ không có trong
văn bản (xem Mục~\ref{sec:invariants}).

\subsection{Mục tiêu}
\label{sec:goals}

Đề tài T11 trong danh sách đề tài của môn học yêu cầu: \emph{TF-IDF retrieval baseline;
multilingual BERT, PhoBERT encoder + QA head hoặc XLM-R; phân tích ảnh hưởng của độ dài
context và loại câu hỏi; đánh giá bằng Exact Match và token-level F1; phân tích câu hỏi
single-sentence so với multi-sentence reasoning.} Bảng~\ref{tab:goals} ánh xạ từng yêu
cầu với nội dung thực hiện và vị trí trình bày trong báo cáo.

\begin{table}[htbp]
\centering
\caption{Đối chiếu yêu cầu đề tài T11 với nội dung thực hiện}
\label{tab:goals}
\renewcommand{\arraystretch}{1.18}
\small
\begin{tabularx}{\textwidth}{@{}L L c l@{}}
\toprule
\textbf{Yêu cầu của đề tài} & \textbf{Nội dung thực hiện} & \textbf{Trạng thái} & \textbf{Mục} \\
\midrule
TF-IDF retrieval baseline & Truy hồi câu bằng TF-IDF (n-gram 1--2) + cosine & \yes & \ref{sec:baseline} \\
Multilingual BERT + QA head & Fine-tune \code{bert-base-multilingual-cased} trên toàn bộ train & \yes & \ref{sec:finetune} \\
XLM-R & \code{deepset/}\allowbreak\code{xlm-roberta-base-squad2}, zero-shot & \yes & \ref{sec:models} \\
PhoBERT encoder + QA head & Không khả thi về kỹ thuật (thiếu fast tokenizer); thay bằng ViSoBERT, encoder tiếng Việt của UIT & \textcolor{warnorange}{thay thế} & \ref{sec:phobert} \\
Exact Match, token-level F1 & Tự cài theo quy ước SQuAD~2.0, điều chỉnh cho tiếng Việt & \yes & \ref{sec:metric-design} \\
Ảnh hưởng độ dài context & Phân tích theo 4 nhóm độ dài, kèm cờ nhóm nhỏ & \yes & \ref{sec:by-length} \\
Single- vs multi-sentence & Heuristic gán nhãn phạm vi suy luận & \yes & \ref{sec:by-qtype} \\
Nêu rõ giới hạn dữ liệu & Mục riêng về giới hạn; khoảng tin cậy cho các số chính & \yes & \ref{sec:limitations} \\
\bottomrule
\end{tabularx}
\end{table}

Ngoài các yêu cầu bắt buộc, nhóm đặt thêm ba mục tiêu về chất lượng quy trình:

\begin{enumerate}
  \item \textbf{Tái lập được.} Mọi số liệu được sinh bởi mã nguồn, ghi ra tệp JSON kèm
        commit, thời điểm, thiết bị, tên split và cỡ mẫu $n$.
  \item \textbf{Chống tự lừa mình.} Đăng ký giả thuyết định lượng \emph{trước} khi chạy
        và định nghĩa trước các ``tín hiệu báo động'' cần điều tra thay vì ăn mừng.
  \item \textbf{Sản phẩm chạy được.} Ứng dụng web cho phép nhập đoạn văn và câu hỏi
        bất kỳ, chọn mô hình và điều chỉnh ngưỡng ``không có đáp án''.
\end{enumerate}

\subsection{Phạm vi}

Đồ án giới hạn trong bài toán hỏi đáp \emph{trích xuất trên một đoạn văn cho trước}.
Các hướng sau nằm \textbf{ngoài phạm vi}: hỏi đáp sinh văn bản (generative QA), truy hồi
đoạn văn từ kho tài liệu lớn (open-domain QA, RAG), hỏi đáp nhiều lượt hội thoại, và
đánh giá trên dữ liệu ngoài miền Wikipedia. Mọi kết luận trong báo cáo chỉ có giá trị
trong phạm vi tập validation của UIT-ViQuAD~2.0 và cấu hình được mô tả.

\subsection{Đóng góp chính}

\begin{enumerate}
  \item Một pipeline MRC tiếng Việt hoàn chỉnh, được tổ chức thành các module có kiểm
        thử (474 test), từ nạp dữ liệu đến ứng dụng web.
  \item So sánh có kiểm soát bốn mô hình, thiết kế để mỗi bước cô lập một biến: so
        khớp từ khoá, chuyển giao đa ngữ zero-shot, fine-tune trong miền, và pretraining
        chuyên tiếng Việt.
  \item Áp dụng nhất quán, ở mọi bảng kết quả, cách tách \emph{câu có đáp án} và \emph{câu
        không có đáp án} --- quy ước đã có sẵn trong công cụ chấm chính thức của SQuAD~2.0,
        chứ không phải một phép đo do nhóm phát minh --- để cho thấy điểm tổng che khuất hai
        kỹ năng tách biệt (tìm đúng biên đáp án và biết từ chối) mà các mô hình mạnh, yếu rất
        khác nhau. Đóng góp là ở thói quen báo cáo, không phải ở bản thân phép tách.
  \item Phân tích nguyên nhân thất bại của ViSoBERT theo quy trình loại trừ, gắn với
        bằng chứng đo được về mức độ phân mảnh token.
  \item Phát hiện và xử lý bốn lỗi ``sai âm thầm'' có thể làm sai kết quả mà không báo
        lỗi: test split không có nhãn, giới hạn số cửa sổ của thư viện
        \code{transformers}, mẫu đánh giá thiên lệch chủ đề và ngưỡng độ dài đáp án
        phụ thuộc tokenizer.
\end{enumerate}

\subsection{Cấu trúc báo cáo}

Phần còn lại của báo cáo gồm: \textbf{Mục~\ref{sec:theory}} trình bày cơ sở lý thuyết
về MRC, TF-IDF, kiến trúc Transformer, các mô hình tiền huấn luyện và thang đo.
\textbf{Mục~\ref{sec:method}} mô tả kiến trúc hệ thống và các quyết định thiết kế.
\textbf{Mục~\ref{sec:experiments}} trình bày môi trường, dữ liệu, kết quả và phân tích.
\textbf{Mục~\ref{sec:conclusion}} kết luận, nêu giới hạn và hướng phát triển.
Phụ lục cung cấp hướng dẫn tái lập và cấu trúc mã nguồn.
